\subsection*{S8. Empty object}

A field whose value is an object with no members, at every document in the
corpus. If some documents supply members and others do not, ordinary inference
applies and the members are simply partial under S4; this scenario is only the
case where the object is empty everywhere.

\begin{lstlisting}[style=json]
{ "a": 1, "b": {} }
\end{lstlisting}

\options
\begin{enumerate}[label=\textbf{(\alph*)}]
  \item \textbf{A table with only an identifier column}, and a module whose
        row type carries nothing but that identifier. \selected
  \item \textbf{Reject}, with an error naming the path.
  \item \textbf{Drop the field} from the schema.
  \item \textbf{\texttt{b : unit}} --- a marker in the parent, with no table
        and no module.
\end{enumerate}

\decision{Emit a table with only an identifier column. The parent holds a
foreign key to it exactly as in S7.}

\begin{lstlisting}[style=ml]
module B : sig
  type id
  type t = { id : id }
  val get : db -> id -> t
end

module Root : sig
  type id
  type t = { id : id; a : int; b_id : B.id }
  val get : db -> id -> t
  val b : db -> t -> B.t
end
\end{lstlisting}

\rationale{The same principle as S6: keep the field, and let its generated
form state exactly what is known about it. Option (a) also preserves a
distinction the others lose. \texttt{\{"b": \{\}\}} and \texttt{\{\}} are
different inputs: the first yields a non-null foreign key to an empty row,
while the second makes \texttt{b} partial under S4. Under (c) and (d) that
difference disappears.

Option (d) is the tempting shortcut, and it holds only until some document in
the corpus gives \texttt{b} a member --- at which point the representation has
to change from \texttt{unit} to a module, rather than the existing table
merely gaining a column.}

\limitation{The table carries no data. Like the \texttt{unit option} of S6, it
exists to record structure that was present in the input, and every row in it
is a bare generated identifier.}
